\chapter{Introduction}
\label{chap:introduction}

\section{Motivation}
Execution order determines when a program accesses its data. A compiler can
therefore change the opportunities for cache reuse by transforming a loop nest,
provided that the transformed program preserves the required semantics.
\citet{wolfLam1991locality} present a locality-oriented loop transformation
algorithm; \citet{mckinleyCarrTseng1996} provide another foundation for examining
locality through loop transformations.

\exampleparagraph{The citations above illustrate narrative citations using
\texttt{\string\citet}. A parenthetical citation can instead be written using
\texttt{\string\citep}, as in the next paragraph.}

The selected literature gives a starting point for discussing locality alongside
other optimisation concerns, including instruction scheduling and memory-operation
costs \citep{wolfMaydanChen1996,carr1996,carrKennedy1994}.

\section{Scope and research questions}
The report can be organised around the following questions:
\begin{enumerate}
  \item How do the selected papers represent data reuse and locality?
  \item Which loop transformations are considered, and how is legality checked?
  \item How do the approaches estimate profitability and account for the target machine?
  \item What evidence supports their conclusions, and what are their limitations?
\end{enumerate}
\todo{Adapt these questions to the papers you have read and the assignment brief.}

\section{Report organisation}
\Cref{chap:background} introduces notation and a simple cost model.
\Cref{chap:literature} provides a framework for comparing the papers.
\Cref{chap:example} demonstrates a loop transformation, and
\cref{chap:evaluation} outlines an evaluation methodology.
\Cref{chap:conclusion} is reserved for the final synthesis.
